\section{Discussion}
\label{sec:discussion}

\subsection{What the detector learned}

The two cross-generator results are not in tension once the variable that
actually changed is identified. Withholding Suno or Udio changes which
generative process produced the synthetic class; the real class, its corpus, its
provenance and its processing chain all stay fixed, and the detector transfers
almost perfectly. Moving to FakeMusicCaps changes the generative processes
\emph{and} replaces the real class with recordings from a different corpus, and
the detector collapses---specifically, by calling every genuine recording
synthetic.

A detector that had learned the artefacts of generation would degrade on unseen
generators and remain calm on unseen real music. Ours does the opposite. The
most economical reading is that it has learned a model of what real music looks
like in its training corpus, and treats departures from that model as evidence
of synthesis. Under that reading a new generator is simply another departure and
is handled correctly, while genuine music from an unfamiliar corpus is also a
departure and is handled catastrophically.

This has a practical consequence for deployment. A detector validated on
cross-generator splits alone can appear robust and still be unusable, because
the population of real music it will meet in production is never the population
it was trained on. The failure it will exhibit is the expensive one: false
accusation.

\subsection{Why $F_1$ is the wrong headline}

Our cross-corpus $F_1$ of 0.750 exceeds the published cross-generator baseline of
0.629, and means nothing---it is exactly the score of a classifier that labels
everything synthetic. The metric is not defective; it is being asked to do
something it cannot, which is to distinguish a detector from a constant
function when classes are unbalanced.

Two cheap conventions prevent this. Report the all-positive baseline beside every
$F_1$, and report false-positive rate alongside it. Either would have caught our
degenerate result immediately. We would encourage their adoption as a norm in
this literature, where class balance varies widely between evaluation sets and
the cost of a false positive is the main reason the work is being done.

\subsection{Calibration as a first-class result}

The robustness table shows AUROC holding above 0.989 under four of six
degradations while $F_1$ falls by up to 0.233, and the cross-corpus evaluation
shows calibration error rising from $1.5\times10^{-4}$ to 0.365. In both cases
the ranking survives and the decision rule does not.

The mechanism is visible in threshold selection. Where validation separates the
classes perfectly, the usual false-positive-budget threshold lands against the
real class and any distribution shift sweeps the entire real population across
it. Placing the threshold at the midpoint of the separating margin is a one-line
change that moved false-positive rate from 1.00 to 0.00 in our cross-generator
evaluation. Detectors intended for forensic use should report the threshold rule
as part of the method, not as an implementation detail.

\subsection{Explanation, and what it is evidence of}

The faithfulness asymmetry---every explanation of real audio beating its random
baseline, none of generated audio doing so---is a negative result about the
method, not about the model. It says that whatever the detector uses to identify
synthesis is not concentrated in time, and therefore that temporal attribution,
which dominates the audio-explainability literature, cannot surface it.

That is consistent with the frequency-domain characterisation of AI-music
artefacts~\cite{afchar2025ismir}, and it is why we separated the frequency claim
from the temporal one in Section~\ref{sec:method} rather than presenting a
single Grad-CAM-style saliency map over a spectrogram the model never sees.

The layer-level comparison carries a second warning. An attention-pooling
architecture hands the analyst a ready-made explanation---the softmax over
layers---and it is tempting to report it, because it costs nothing and looks
like an answer. In our model it is flat to within 1.09$\times$ and its ordering
is negatively rank-correlated with the exact Shapley attribution over the same
layers in every audited song. The attention weights are not a weak version of
the Shapley values; they point the other way. Architectural attention is
evidence about the architecture, not about the decision, and papers that report
it as an explanation should be read accordingly. This is the concrete payoff of
running Levels~1 and~3 as a contrast rather than as two independent views: the
disagreement is the finding.
